\chapter{Alcohol, Aroma, and Sourness}

\begin{kaichang}
Imagine three things in front of you: white vinegar, a fully ripe banana, and a small bottle of medicinal disinfectant alcohol.

Before calling them unrelated, try some thought experiments. Open wine left for weeks often becomes sour with a hint of vinegar; overripe bananas give off a faint alcoholic scent; fermenting grains or fruit can produce either alcohol or vinegar. An invisible production line seems to connect them, transforming one substance into another.

Guess what that line is. How can alcohol, fragrance, and sourness, such different experiences, be three stations along it? By the end, you will draw the whole line using just four components.
\end{kaichang}

\section{Three Liquids, One Invisible Production Line}

First stay with what eyes, nose, and skin report directly.

White vinegar is a clear, colorless liquid with a sharp sour smell up close. A little on the tongue makes you squint. It mixes with water in any proportion, never separating.

Medicinal alcohol is also clear and colorless, but behaves differently: on the back of a hand it disappears within seconds, leaving coolness as evaporation removes heat, discussed in Chapter 4. It too mixes freely with water, and ignites easily with a quiet pale-blue flame.

Banana aroma is invisible and intangible, yet reaches your nose on its own. Some molecules must continually escape the flesh, drift through air, and land on nasal receptors. Anything smelled must involve readily evaporating molecules.

Three substances, three personalities, yet one quiet commonality: their molecular formulas contain carbon, hydrogen, and oxygen, with often only one or two oxygen atoms. How can so few oxygens have such a presence?

\section{Oxygen Atoms: Interfaces on Carbon Frameworks}

Chapter 36 made carbon a master framework builder: four bonds, chains, and rings. Pure hydrocarbon frameworks such as Chapter 39's alkanes, however, are dull, water-insoluble, and unreactive, with burning their most exciting moment.

Insert oxygen and everything changes. A famous electron magnet, as Chapter 11's electronegativity showed, it pulls nearby electrons toward itself, creating locally charged hotspots in a previously even framework. Uneven electron distribution invites water to dock and other molecules to attack. Chapter 37 called these sites \textbf{functional groups}, molecular reaction interfaces.

Only four interfaces star here. Meet them first:

\begin{itemize}[itemsep=2pt]
\item \textbf{Hydroxyl}: oxygen joined to hydrogen, $-\mathrm{OH}$, attached to a carbon chain makes an \textbf{alcohol}. Ethanol and methanol are examples.
\item \textbf{Carbonyl}: carbon double-bonded to oxygen, $\mathrm{C{=}O}$. At a chain end it gives an \textbf{aldehyde}; in the middle, a \textbf{ketone}. Ethanal is an aldehyde; acetone in remover is a ketone.
\item \textbf{Carboxyl}: carbonyl and hydroxyl on the same carbon, $-\mathrm{COOH}$, acquire a new name and character as a \textbf{carboxylic acid}. Vinegar's acetic acid belongs here.
\item \textbf{Ester group}: replace carboxyl's hydrogen with a carbon chain, $-\mathrm{COO}-$, and obtain an \textbf{ester}. Many banana and pineapple aroma molecules are esters.
\end{itemize}

These are four variations on one theme, combining carbonyl and hydroxyl components. The invisible production line is oxygen changing outfits on carbon frameworks. Follow it step by step.

\section{Hydrogen Bonds: Solubility's Handshakes}

First solve an observable puzzle: why do alcohol and vinegar mix freely with water while Chapter 39's gasoline stubbornly separates?

Chapters 22 and 23 described water's specialty, hydrogen bonding. Water molecules pull one another into a sticky network. To dissolve, an outsider must communicate with that network by forming hydrogen bonds too; otherwise water excludes it into a separate layer.

Hydroxyl is hydrogen bonding's standard interface: slightly negative oxygen and slightly positive hydrogen fit water perfectly. Ethanol joining water resembles one water droplet joining another. Acetic acid's carboxyl also shakes hands and permits unrestricted mixing. Gasoline's alkanes lack any such interface, and the water network shuts them out.

\begin{qiaoliang}
Hydroxyl lets alcohol shake hands, but the carbon chain remains unsociable. Solubility becomes a tug-of-war: hydroxyl pulls toward water, the chain toward oil. Size decides. Methanol, \chem{CH_3OH}, and ethanol, \chem{CH_3CH_2OH}, have short chains, so hydroxyl wins and they mix freely. Four-carbon butanol dissolves only about eight grams per hundred grams of water at $20\,^{\circ}\mathrm{C}$. Eight-carbon octanol dissolves about 0.05 grams per hundred grams: the long chain drowns out hydroxyl. The same rule explains boiling points. Ethanol boils at $78\,^{\circ}\mathrm{C}$; remove hydroxyl, leaving ethane's hydrocarbon framework, and it boils at $-89\,^{\circ}\mathrm{C}$, nearly 170 degrees lower. That difference prices the hydrogen-bond network: boiling ethanol requires extra energy to pull apart all those molecular handshakes.
\end{qiaoliang}

\begin{shiyan}
\textbf{Which Liquid Gets Along with Water?}

Materials: three clear glasses; clean water; a little cooking oil; white vinegar; a little spirits or medicinal alcohol; a chopstick.

Procedure: half-fill all glasses with water. Add a few oil drops to the first, a spoonful of vinegar to the second, and a spoonful of alcohol to the third. Stir each ten times, stand for five minutes, and view from the side.

What to observe: which has distinct layers? Which is uniform? Held against light, are uniform samples truly clear or faintly cloudy? Predict replacing alcohol with equal gasoline; Chapter 39's readers should know.

Safety: use small amounts only and never taste experimental liquids. Keep alcohol away from flames and stoves. Afterward, pour liquids down the drain and flush with clean water.
\end{shiyan}

\section{Alcohol to Vinegar: An Oxidation Chain}

Now enter symbols and draw the opening line. Yeast converts sugar to ethanol, unpacked in Chapter 52, but ethanol is unstable in the presence of oxygen and certain bacteria, oxidizing stepwise:

\[\chem{CH_3CH_2OH} \rightarrow \chem{CH_3CHO} \rightarrow \chem{CH_3COOH}\]

Ethanol, ethanal, acetic acid. Count oxygens: one, one, two. Track groups: hydroxyl becomes aldehyde carbonyl, then gains oxygen to become carboxyl. Wine turning sour means airborne acetic-acid bacteria have completed this chain. Kitchen vinegar and wine really are two stations apart.

Chapter 33 defined oxidation as losing electrons, not necessarily meeting oxygen gas. Each step means carbon's electrons are partly taken by oxygen and its oxidation state rises. Bodily enzymes and vinegar bacteria perform the same electron transport, one in cells and one in a vinegar vat.

Meet a dangerous relative: methanol, \chem{CH_3OH}, with one fewer carbon. It follows the same oxidation chain but gives formaldehyde, \chem{HCHO}, and formic acid, \chem{HCOOH}. Formaldehyde damages cellular proteins, while accumulated formic acid disturbs blood's acid--base balance (Chapter 32), with optic nerves among the first victims. Industrial alcohol often contains methanol: this is counterfeit liquor's chemical route to blindness, not exceptional strength but a different oxidation branch leading to toxins.

\section{After Alcohol Enters the Body}

Ethanol is small, water-friendly, and not too averse to fat because of its short two-carbon oily tail. It therefore travels almost freely through the body, entering blood from stomach and small intestine, crossing membranes, and reaching the brain within minutes. There it interferes with molecular switches between nerve cells, strengthening some slowing signals and weakening accelerating ones, beginning slower reactions and unsteady walking. That is physiology, not our focus; follow the chemistry of dismantling this foreign molecule.

First, solve a disinfectant puzzle: why is pharmacy alcohol about 75\%, not pure 100\%? Can stronger concentration weaken germ killing? In a sense, yes. Alcohol enters microbes and denatures proteins. Pure alcohol acts too hastily, immediately coagulating surface proteins into a shell that impedes entry. Dilution with water allows more gradual penetration before internal action. Remember this challenge to stronger-is-better and purer-is-better intuition: ratios often matter more than purity.

The liver is the main dismantling site, using the same oxidation chain but two enzymes rather than vinegar bacteria. Chapter 49 explains enzymes:

\begin{itemize}[itemsep=2pt]
\item Alcohol dehydrogenase oxidizes ethanol to ethanal, also called acetaldehyde;
\item Aldehyde dehydrogenase oxidizes acetaldehyde to acetic acid;
\item Acetic acid enters whole-body energy metabolism, breaking into carbon dioxide and water and releasing energy (Chapter 53).
\end{itemize}

The conspicuous intermediate is acetaldehyde, much more toxic than ethanol. It dilates blood vessels, reddening face and neck, causes headache, palpitations, and nausea, and attacks DNA directly. The International Agency for Research on Cancer places alcoholic beverages themselves and acetaldehyde associated with drinking in Group 1, sufficient evidence of human carcinogenicity. Understand this precisely: it describes evidence strength, not a claim that one drop causes cancer. Risk increases with dose, toxicology's first firm rule and Chapter 1's dose principle.

The second firm rule is \textbf{individual variation}. Everyone has enough first-step enzyme; the bottleneck is the second. About one-third of East Asians carry markedly inefficient aldehyde dehydrogenase, explored in the next evidence story. Acetaldehyde remains longer and reaches higher concentrations, producing much greater exposure from the same drink. Thus \textbf{no safe dose is identical for everyone}; for minors with developing brains, alcohol interference has no claim to safety at all. This chapter explains chemistry and evidence, offering no drinking recommendations.

\begin{qiaoliang}
Estimate one 330-milliliter can of 5\% beer: $330\times 5\% \approx 16.5$ milliliters of ethanol. At \ud{0.79}{g/mL}, that is about thirteen grams; divided by \ud{46}{g/mol}, about 0.28 moles or $1.7\times 10^{23}$ molecules. An adult liver processes about seven grams hourly on average, with severalfold individual differences. At that speed, this one can requires about two hours of continuous liver work. Calculate how molecules accumulate if intake increases or accelerates. Drunkenness essentially means entry far outrunning breakdown.
\end{qiaoliang}

\begin{wuqu}
``People who flush can hold their drink; practice makes them better.'' Exactly backward. Flushing warns of acetaldehyde accumulation and forced vasodilation: the second breakdown step is lagging and toxic material remains. Long-term drinkers with this trait face significantly greater risks of esophageal cancer and other disease than nonflushing drinkers. Acquired tolerance merely dulls the brain to interference, turning the alarm down while acetaldehyde's attack on DNA continues. A silent alarm does not mean the fire is out.
\end{wuqu}

\section{How Scientists Solved the Flushing Mystery}

\begin{zhengju}
Alcohol flushing is common in East Asians and rare in Europeans. An everyday curiosity gave scientists a standard question: where exactly does the difference arise?

Candidates included diet, vaguely defined constitution, or psychology. In the 1970s, researchers converted the question into measurable chemistry: in the ethanol--acetaldehyde--acetic acid chain, which link slows in people who flush? Comparing liver dehydrogenase activities revealed a pattern: many flushers had entirely normal alcohol dehydrogenase but slow aldehyde dehydrogenase. The first enzyme rapidly made acetaldehyde, while the second lagged, causing accumulation. This directly predicted higher blood acetaldehyde, confirmed by measurement.

Why was the enzyme slow? By the late 1980s, DNA provided an answer: one base difference in its gene replaced amino acid 487, altering the three-dimensional protein and leaving its assembled machine loose and inefficient. Later large epidemiological studies compared genotype, drinking, and esophageal-cancer incidence, confirming an unsettling association: carriers who continued drinking long-term had multiply increased risk.

What makes this story elegant? Every step excludes alternatives. Not mystical constitution, because specific enzyme activities are measurable; not drinking speed, because genotype predicts acetaldehyde under equal intake. The final chain---one base, one amino acid, a distorted enzyme, accumulated acetaldehyde, greater risk---has independent evidence at each link. It previews Volume III's central drama: tiny DNA differences amplify through protein shape into visible individual variation, fully developed in Chapters 71 and 72.
\end{zhengju}

\section{Carbonyl: An Electron-Pulling Double Bond}

Aldehydes and ketones are the line's intermediate station, sharing carbonyl's carbon--oxygen double bond. Chapter 38's general approach starts where electrons crowd and where they are scarce. Carbonyl is textbook: oxygen draws electrons strongly toward itself, becoming slightly negative and leaving a slightly positive electron gap at carbon.

That gap is the entire source of carbonyl reactivity. Electron-rich molecules, nucleophiles, are attracted and form new bonds, making carbonyl carbon among organic chemistry's most approachable sites. Your body exploits this extensively: glucose too has a carbonyl, and cellular energy factories begin glucose breakdown by manipulating it (Chapters 46 and 52).

Aldehydes and ketones differ only in position. An aldehyde's terminal carbonyl also bears hydrogen, easily taken by oxygen, so it readily oxidizes to a carboxylic acid, as acetaldehyde becomes acetic acid. Ketones have carbon on both sides and no convenient hydrogen, so are steadier. Classic laboratory differentiation uses a mild oxidizer that attacks aldehydes but not ketones: the oxidized one is the aldehyde.

Everyday carbonyls abound: rapidly evaporating remover smells of acetone, the simplest ketone and an excellent solvent for plastic and nail-polish polymers. Formalin preserving biological specimens is aqueous formaldehyde; troubling odors in newly decorated rooms often include such small molecules.

\section{Why Acids Are Sour and Esters Fragrant}

The third station is carboxylic acid. Chapter 32 defined acids as molecules willing to release hydrogen ions, \chem{H^+}. But alcohol hydroxyl also has hydrogen: why is ethanol not acidic while acetic acid plainly is?

Ask how life feels after letting go. A proton donor's remainder must bear a negative charge. The more crowded and homeless that charge, the less willing the molecule is to release hydrogen. Ethanol leaves it wholly on one oxygen, so rarely releases hydrogen and is neutral. Acetic acid shares it across \textbf{two oxygens in carboxylate}, Chapter 41's delocalization, with electrons not confined to one atom. Halving the burden makes release easier:

\[\chem{CH_3COOH} \rightleftharpoons \chem{CH_3COO^-} + \chem{H^+}\]

The arrow is reversible: weak acetic acid has only a few willing donors per hundred molecules, the rest hesitating in Chapters 31 and 32's equilibrium. Yet this acidity makes taste buds cry out, bubbles baking soda, and dissolves scale. Vinegar's kitchen magic comes from those few hydrogen ions.

Bring two characters together: carboxylic acid and alcohol contribute from carboxyl and hydroxyl, expel water, and join the remainder into ester:

\[\chem{CH_3COOH} + \chem{CH_3CH_2OH} \rightleftharpoons \chem{CH_3COOCH_2CH_3} + \chem{H_2O}\]

This is \textbf{esterification}, yielding sweet, fruity ethyl acetate. Reversible again, it proceeds sluggishly at room temperature with both sides stalemated. Industry uses acid catalysis for speed, not a changed endpoint (Chapter 30), and removes produced water to push equilibrium toward ester.

Esters lead fruit aromas: ripe banana mainly owes sweetness to isoamyl acetate, pineapple includes ethyl butyrate, and apples and pears have their own ester recipes. Think of evolution: unripe fruit is sour and astringent, refusing to be eaten before seeds develop. Mature seeds bring more sugar, less acid, and abundant ester synthesis, a ready-to-eat advertisement recruiting animals to disperse them. Banana scent is advertising plants have written for hundreds of millions of years; Chapter 76 explains natural selection's authorship.

Aroma esters are usually small and poor at handshaking with water, so readily evaporate into noses. Conversely, water may slowly hydrolyze esters back into acids and alcohols; bases greatly accelerate and complete this. That property leads to two everyday products.

\section{From Soap to Biodiesel}

Chapter 47 formally introduces fats: three fatty acids linked by esters to glycerol. If esters hydrolyze, so do fats. Boil fats with caustic soda, \chem{NaOH}, cutting ester bonds, releasing glycerol, and making sodium fatty-acid salts: \textbf{soap}. Its special reaction name is saponification.

We already prepared soap's cleaning explanation. Sodium fatty-acid salts have a long, unsociable, oil-loving chain at one end and a charged, water-loving carboxylate at the other. Chains embrace grease while water pulls the charged ends, dragging whole grease clumps into water. One molecule belongs to both camps: a beautiful application of Chapter 25's like-dissolves-like rule.

\textbf{Translator safety note:} The following biodiesel paragraph describes industrial chemistry, not a home recipe. Do not handle methanol or caustic catalysts for this activity, or put an experimental product into an engine. Consult \href{https://institute.acs.org/acs-center/lab-safety/education-training/safer-experiments.html}{ACS guidance on chemical demonstration safety}.

Ester chemistry also reaches fuel stations. Mix vegetable or used cooking oil with methanol and catalyze transesterification, replacing glycerol at ester groups with small methyl groups. Fatty-acid methyl esters result, burnable directly as diesel: \textbf{biodiesel}. Glycerol byproduct can go to personal-care manufacturers. A vehicle fueled from discarded frying oil sounds like a joke, but real transesterification lies behind it: chemistry recognizes structure, not prestigious origins.

\section{When Does the Group Manual Fail?}

You may feel you have a universal manual: identify the group, predict behavior. Time for cold water: the manual has disclaimers.

First, neighbors rewrite group behavior. Hydroxyl on ethanol's chain hardly releases hydrogen, but within acetic acid's carboxyl it does, because adjacent carbonyl shares negative charge. Molecules are not merely assembled parts; parts communicate.

Second, large molecules have many interfaces, and circumstances decide the loudest. Glucose has five hydroxyls and a carbonyl; its behavior is the whole committee's decision, not one interface's view (Chapter 46).

Third, sensory properties such as odor are especially unruly. Similar esters can smell very different, while distant structures sometimes smell alike. Olfactory receptors recognize three-dimensional shape puzzles, as Chapter 42 will stress, not lists of groups.

Fourth, dose and individual differences always matter. Discussing harm without dose, or assuming identical responses from everyone, mistakes the manual for law.

\section{Back to the Three Liquids}

Now complete the promised line:

\begin{itemize}[itemsep=2pt]
\item \textbf{Medicinal alcohol} is ethanol. Hydroxyl makes it water-miscible, volatile, and cooling; in the body it follows ethanol--acetaldehyde--acetic acid breakdown.
\item \textbf{White vinegar} is dilute acetic acid. Carboxyl willingly releases hydrogen ions, producing unmistakable sourness. Open wine souring means acetic-acid bacteria completing the chain.
\item \textbf{Banana's sweet scent} mainly comes from esters, fruit acids and alcohols joining while expelling water. Overripe banana's alcohol smell comes from oxygen-starved flesh cells fermenting sugar into ethanol.
\end{itemize}

They occupy not three separate streets but the intersection of oxidation and esterification: alcohol, aldehyde or ketone, carboxylic acid, ester, four interfaces with routes between them. More importantly, everything upstream points to \textbf{sugar}. Yeast makes alcohol from sugar; bacteria make vinegar from alcohol; fruit builds acids and esters from sugar to attract animals. Where does sugar come from, and why is it both fuel and material? That is the next part: Chapter 46 introduces one of life's busiest molecules.

\begin{tiaozhan}
Acid strength has a precise number, $\mathrm{p}K_a$, the negative logarithm of its ionization equilibrium constant (Chapters 31 and 32). Smaller means stronger. Acetic acid's $\mathrm{p}K_a$ is about 4.8; ethanol's about 16. An eleven-unit logarithmic gap means roughly $10^{11}$ greater hydrogen-donating tendency: ten billionfold! The same hydroxyl hydrogen, with just an adjacent carbonyl sharing negative charge, differs in value by ten billionfold. For why delocalization spreads charge and how resonance drawings work, Chapter 41 starts from benzene. Skipping this box leaves the main thread intact.
\end{tiaozhan}

\section*{Try It Yourself}

\begin{sikaoti}
\item Diagram the ethanol--acetaldehyde--acetic acid oxidation chain, using boxes with hydroxyl, carbonyl, and carboxyl labels. Beside each arrow, indicate whether oxygen was added or hydrogen lost. What same work do bodily enzymes and vinegar bacteria perform?
\item Compare three-carbon 1-propanol and glycerol, with one hydroxyl on each of three carbons. Which dissolves more readily and evaporates less? Explain with the hydroxyl--chain tug-of-war and predict boiling-point order.
\item A little acetic acid and baking soda, sodium bicarbonate, bubble together. Write acetic acid's ionization equation and identify the gas. Does the experiment prove a strong or weak acid? Why?
\item Draw a soap material-flow diagram: a long molecule labeled oil-loving and water-loving at opposite ends, grease, and water molecules, with arrows showing grease movement after stirring. Why can plain and soapy water both remove dust, but only soap tackles grease?
\item Ethyl acetate hydrolyzes back to acid and alcohol. Using Chapter 31's equilibrium, why does a factory remove water continuously to make more ester? Can a catalyst increase the final ester yield?
\item Two classmates drink equal alcohol; one flushes, one does not. Explain the molecular reasons using the metabolic chain, and why saying flushing means better tolerance reverses causality.
\end{sikaoti}
